\documentclass[12pt]{article}
\usepackage[margin=1in]{geometry}
\usepackage{enumitem}
\usepackage{titlesec}
\usepackage{parskip}
\usepackage{amsmath}
\usepackage{hyperref}
\usepackage{fontenc}

\title{\textbf{Boix 2003 Claude Guide}\\
\large Boix, Carles. \textit{Democracy and Redistribution}.\\
New York: Cambridge University Press, 2003.}
\author{}
\date{}

\begin{document}

\maketitle
\tableofcontents
\newpage

%--------------------------------------------------
\section{Central Claims}
%--------------------------------------------------

\begin{itemize}[leftmargin=*, itemsep=4pt]
    \item Boix offers a formal, unified theory of \textbf{democratization and authoritarian regime persistence} grounded in the redistributive conflict between the rich (or the asset-holding elite) and the poor (the median voter/majority). Whether a country democratizes, remains authoritarian, or reverts from democracy to dictatorship is explained by the interaction of two structural variables: the level of \textbf{economic equality} and the degree of \textbf{capital mobility} (the specificity or fungibility of the assets the elite hold).
    \item Democracy is, at its core, a redistributive institution: extending the franchise shifts the decisive voter from a wealthy elite to a poorer median citizen, who will use the state's fiscal and regulatory power to redistribute income and wealth downward. Elites who anticipate large redistributive losses under democracy will resist democratization; elites who anticipate small losses will tolerate or even initiate it.
    \item The key insight distinguishing Boix from a simple Meltzer-Richard extension is that the \textbf{cost of democratization to elites depends not only on inequality but on asset mobility}. When wealth is held in fixed, immobile assets (land), it is easy for a democratic majority to identify, tax, and expropriate; when wealth is mobile (financial capital, human capital, portable industrial capital), elites can threaten capital flight, disinvestment, or relocation, which disciplines the redistributive appetite of the median voter and lowers the effective cost of democracy to elites.
    \item This produces Boix's central comparative-statics result: \textbf{democratization is most likely when equality is high and/or asset mobility is high}; \textbf{authoritarianism (and elite repression of democratic movements) is most likely when inequality is high and assets are immobile (land-based)}---the paradigmatic case of landed oligarchies violently resisting the extension of suffrage.
    \item Democratic \textbf{survival/breakdown} is governed by the same logic in reverse: democracies are more likely to break down (via elite-sponsored coups) when inequality is high and capital is immobile, because the poor majority's ability to expropriate the rich under democratic rules becomes intolerable to elites who cannot exit with mobile assets.
    \item The decision to repress or tolerate is modeled as a strategic choice under uncertainty: elites compare the (probabilistic, risky, and costly) option of repression against the expected redistributive tax burden of democracy. Repression itself has costs (military/organizational costs, risk of civil conflict, potential loss of legitimacy and investment), so as repression becomes more costly or less likely to succeed, elites become more willing to tolerate democratic transition even holding inequality fixed.
    \item Boix explicitly rejects purely culturalist (civic-values), purely modernization-theoretic (development-as-automatic-cause), and purely institutionalist accounts of democratization in favor of a materialist, distributive-conflict account rooted in the economic interests of contending classes---closely paralleling but formally extending the class-coalitional logic of Barrington Moore.
\end{itemize}

%--------------------------------------------------
\section{Core Theoretical Model}
%--------------------------------------------------

\begin{itemize}[leftmargin=*, itemsep=4pt]
    \item \textbf{Baseline redistribution game (Meltzer-Richard logic)}: Under universal suffrage, the median voter is decisive. If the median income is below the mean (as it is in essentially all real-world distributions), the median voter prefers a positive tax rate financed redistribution; the more skewed the distribution (the further the median falls below the mean), the higher the equilibrium tax rate the median voter chooses. Boix takes this median-voter redistribution result as the baseline expectation for what democracy \emph{does} once established, and then asks under what conditions elites will allow that regime to come into being at all.
    \item \textbf{The elite's decision problem}: Facing a rising threat of popular mobilization for suffrage extension, the elite (modeled as a unitary actor representing asset-holders) chooses between: (1) \textbf{tolerating} democratization and accepting the resulting redistributive tax rate chosen by the new median voter; or (2) \textbf{repressing} the democratic challenge, at some cost $c$ (which may include forgone investment, international sanction, risk of losing a civil conflict, and the deadweight costs of coercion), and if repression succeeds, retaining authoritarian control (and the associated low or zero redistribution).
    \item \textbf{Two key parameters}:
        \begin{itemize}
            \item \textbf{Equality (or its inverse, inequality)}: determines how large the redistributive tax rate under democracy would be; higher inequality means a democratically-chosen tax rate further from the elite's ideal point, i.e., a larger transfer away from the rich.
            \item \textbf{Asset mobility/specificity}: determines how much of the nominal redistributive burden the elite can actually avoid by moving, hiding, or disinvesting their capital. High mobility acts as a de facto ceiling on the effective tax rate the median voter can impose, because attempting to tax mobile capital too heavily induces flight and destroys the tax base (a self-limiting mechanism akin to the standard capital-mobility/tax-competition literature).
        \end{itemize}
    \item \textbf{Comparative statics}: Elites tolerate democratization when the expected cost of repression exceeds the expected (mobility-discounted) cost of redistribution under democracy. This yields four broad regime configurations: (i) \emph{low inequality + any mobility}---democracy is stable and cheap to tolerate; (ii) \emph{high inequality + high mobility}---democratization is still plausible because mobile capital blunts the effective redistributive threat (Boix's key innovation relative to pure inequality-based accounts); (iii) \emph{high inequality + low mobility (landed elites)}---authoritarianism and repression are likely, since redistribution would be severe and unavoidable; (iv) \emph{repression is very costly regardless of inequality/mobility (e.g., cannot suppress a mobilized population, external pressure)}---democratization is likely even under adverse distributive conditions.
    \item \textbf{Extension to democratic breakdown}: The same repression/tolerate calculus is run by elites \emph{within} an existing democracy: if the realized (or anticipated) redistributive tax rate becomes unbearable relative to the cost of a coup, elites overthrow democracy. This makes regime type endogenous and potentially cyclical, rather than a one-way modernization process---directly contrasting with linear modernization theory.
    \item \textbf{Role of the poor/median voter's own strategic behavior}: Boix's model also allows for the poor to moderate their redistributive demands anticipating the elite's repression/reaction function (a kind of subgame-perfect equilibrium logic), so democratization can be sustained by self-limiting redistribution on the part of the newly enfranchised majority---an equilibrium outcome rather than an assumption.
    \item \textbf{Formal apparatus}: The model is built as a game between two (or more) unitary class actors with quasi-linear utility over post-tax income, standard median-voter tax-choice machinery, and a repression technology with an exogenous cost parameter and success probability---placing it squarely in the rational-choice/formal comparative-historical tradition rather than qualitative-inductive class-coalition analysis.
\end{itemize}

%--------------------------------------------------
\section{Core Empirical Strategy}
%--------------------------------------------------

\begin{itemize}[leftmargin=*, itemsep=4pt]
    \item Boix constructs a \textbf{large cross-national historical dataset} covering regime transitions (democratization and democratic breakdown/reversion to authoritarianism) across a broad sample of countries from the nineteenth century through the late twentieth century, considerably extending the temporal and spatial reach of prior quantitative democratization studies (e.g., Przeworski et al.).
    \item Key independent variables are operationalized as: (1) \textbf{income/wealth inequality} (using available historical inequality proxies, land distribution data, and Gini-type coefficients where obtainable); and (2) \textbf{asset structure/mobility}, proxied by the relative importance of agriculture/land versus industrial and financial capital in the economy (e.g., share of the labor force or capital stock in agriculture versus manufacturing/finance), on the logic that land is the paradigmatic immobile, easily-taxed asset while industrial and especially financial capital is comparatively mobile.
    \item The dependent variable is coded as \textbf{regime transitions}: movement from authoritarian rule to democracy (democratization), and from democracy back to authoritarian rule (breakdown), using standard historical regime classifications (extending/complementing datasets such as Polity and Przeworski et al.'s ACLP/DD dataset).
    \item Boix uses \textbf{event-history/survival analysis and logit/probit models} of regime transition to test whether higher equality and higher asset mobility jointly and interactively predict a higher hazard of democratization and a lower hazard of authoritarian reversion, net of standard controls (income level, British colonial heritage, region, period effects, diffusion/neighborhood democracy).
    \item Historical case illustrations (Britain, Argentina, and other landed-oligarchy versus commercial/financial-elite comparisons) are used alongside the statistical analysis to illustrate the causal mechanism qualitatively---elites in economies dominated by mobile capital (Britain's evolving commercial/financial sector) conceding suffrage relatively peacefully, versus landed oligarchies (parts of Latin America, pre-reform agrarian states) resisting suffrage extension and precipitating repression or civil conflict.
    \item Boix explicitly tests his equality/mobility interaction against rival predictors---most importantly per capita income (the standard modernization-theory proxy)---to show that the distributive/asset-structure variables retain explanatory power (and in his view, superior power) once controls for development level are included.
\end{itemize}

%--------------------------------------------------
\section{Core Works Cited}
%--------------------------------------------------

\begin{itemize}[leftmargin=*, itemsep=4pt]
    \item \textbf{Daron Acemoglu and James Robinson}, \textit{Economic Origins of Dictatorship and Democracy} (2006)---the most direct theoretical sibling and rival: an independently developed but closely related formal model of democratization as a strategic response by elites to the threat of revolution/redistribution from below. Acemoglu-Robinson emphasize democratization as a commitment device to credibly promise future redistribution and avert revolution, whereas Boix emphasizes the ex ante cost-benefit calculus over repression versus toleration given asset mobility; the two frameworks are frequently taught and cited together as the two leading formal political-economy theories of democratic transition.
    \item \textbf{Allan Meltzer and Scott Richard}, ``A Rational Theory of the Size of Government'' (1981, \textit{Journal of Political Economy})---the median-voter model of redistribution that supplies Boix's baseline mechanism for why democratic majorities redistribute, and why greater inequality (median further below mean) produces higher equilibrium taxation.
    \item \textbf{Dietrich Rueschemeyer, Evelyne Huber Stephens, and John D. Stephens}, \textit{Capitalist Development and Democracy} (1992)---a class-coalitional, historical-sociological account of democratization emphasizing the working class and cross-class alliances as the driving force behind suffrage extension; Boix's approach shares the class-conflict starting point but formalizes it as a game-theoretic elite decision problem rather than a coalition-formation/class-power narrative.
    \item \textbf{Seymour Martin Lipset}, ``Some Social Requisites of Democracy'' (1959, \textit{American Political Science Review}) and \textit{Political Man} (1960)---the foundational modernization-theory account that economic development produces the social conditions (literacy, middle class, moderation) conducive to democracy; Boix directly tests and partially rejects this account in favor of the distributive/asset-structure explanation.
    \item \textbf{Barrington Moore Jr.}, \textit{Social Origins of Dictatorship and Democracy} (1966)---the key predecessor whose class-coalitional account of routes to the modern world (bourgeois-led democratic path, landed-elite-led fascist/authoritarian path, peasant-led communist path) Boix can be read as formalizing and generalizing: Moore's ``no bourgeoisie, no democracy'' logic and his emphasis on landed elites as the great obstacle to democratization map directly onto Boix's inequality/asset-mobility parameters.
    \item \textbf{Adam Przeworski, Michael Alvarez, Jose Antonio Cheibub, and Fernando Limongi}, \textit{Democracy and Development} (2000)---the direct empirical and theoretical rival/complement Boix engages most explicitly: Przeworski et al. find that development does not cause democratic \emph{transitions} but strongly predicts democratic \emph{survival} once transitions occur (an ``exogenous'' view of transitions and endogenous view of survival); Boix instead argues that structural economic variables (inequality, asset mobility) causally shape the likelihood of transitions themselves, offering a more structural account than Przeworski et al.'s comparatively agnostic treatment of transition timing.
\end{itemize}

%--------------------------------------------------
\section{Methodological Contributions and Limitations}
%--------------------------------------------------

\begin{itemize}[leftmargin=*, itemsep=4pt]
    \item Boix's book is a leading exemplar of \textbf{formal, rational-choice comparative-historical political economy}: it combines an explicit game-theoretic model with large-N historical hypothesis testing, a combination that was still relatively novel in comparative democratization research at the time and that has become a template for subsequent formal-empirical work (including Acemoglu-Robinson).
    \item \textbf{Data quality concerns}: A recurring critique is that reliable income/wealth inequality data for nineteenth- and early-twentieth-century cases are extremely scarce, and the proxies Boix and others use (land concentration, sectoral employment shares) may not accurately capture the theoretically relevant construct (the effective, politically-salient distribution of income perceived by contemporary elites and citizens). This raises concerns about measurement error large enough to threaten the reported statistical relationships.
    \item \textbf{Endogeneity of asset mobility}: Critics note that asset mobility/specificity is not fully exogenous to politics---capital mobility itself can be a function of prior political institutions, international capital markets, and state capacity to enforce or forgo capital controls, raising the possibility of reverse causation or a jointly determined equilibrium between institutions and asset structure rather than the one-directional causal story the model assumes.
    \item \textbf{Aggregation and unitary-actor assumptions}: Modeling ``the elite'' and ``the poor'' as unitary rational actors abstracts away from intra-class heterogeneity, collective-action problems within each class, and the role of specific political entrepreneurs/organizations---simplifications common to formal models but that limit the theory's ability to explain within-case variation or the precise timing/sequencing of transitions.
    \item \textbf{Overdetermination and comparison with Acemoglu-Robinson}: Because Boix's model and Acemoglu-Robinson's model make similar directional predictions from partially overlapping mechanisms (redistributive threat, elite cost-benefit calculus), it can be difficult to empirically adjudicate between them, and some critics argue the added complexity of asset mobility does not yield sufficiently distinct, falsifiable predictions from a simpler inequality-only account.
    \item \textbf{Strengths}: Despite these critiques, the book is widely credited with formally integrating land/capital structure into democratization theory (a dimension largely absent from Lipset-style modernization accounts and underdeveloped in Moore's more qualitative treatment), and with providing one of the most rigorous large-N tests linking economic structure to regime dynamics across two centuries.
\end{itemize}

%--------------------------------------------------
\section{Connections to the Broader Literature}
%--------------------------------------------------

\begin{itemize}[leftmargin=*, itemsep=4pt]
    \item \textbf{Moore 1966}: The most direct predecessor. Boix's book can be read as taking Moore's qualitative, case-based class-coalition theory of democratization (bourgeois-led democracy, landed-elite-led authoritarianism/fascism, peasant-led revolution) and recasting it as a rigorous game-theoretic model with explicit micro-foundations (elite repression/tolerance calculus) and a large-N cross-national test, trading Moore's rich case narrative for formal generality and statistical leverage.
    \item \textbf{Huntington 1968}: A useful contrast in explanatory strategy rather than a direct rival: Huntington explains political order/decay primarily through the pace of \textbf{institutionalization} relative to mobilization/participation, largely independent of the specific economic-class content of political demands, whereas Boix offers a purely \textbf{economic-structural} account (inequality and asset mobility) of when regime transitions occur, with institutions as an outcome rather than an independent cause.
    \item \textbf{Dahl 1972}: Dahl's \textit{Polyarchy} specifies the descriptive conditions (contestation and inclusiveness) that define the democratic outcome Boix's model is trying to formally explain the emergence of; Boix can be read as supplying a causal, materialist theory of \emph{why} and \emph{when} societies move toward Dahl's polyarchic conditions, rather than a competing definition of democracy itself.
    \item \textbf{Przeworski et al. 2000}: The most direct empirical and theoretical rival/complement. Both works ask what economic conditions produce and sustain democracy, but Przeworski et al. are skeptical that development (or, by extension, most structural economic variables) causes democratic \emph{transitions}---finding development instead robustly predicts democratic \emph{survival}---while Boix argues that inequality and asset mobility structurally determine the timing and likelihood of transitions themselves, offering a more determinate structural theory of transition than Przeworski et al.'s comparatively agnostic stance.
    \item \textbf{Svolik 2012}: \textit{The Politics of Authoritarian Rule} examines power-sharing and repression dynamics \emph{within} authoritarian regimes (the problem of authoritarian elites credibly committing to share power with each other); this is directly complementary to Boix's model of elites' repression-versus-toleration calculus vis-\`a-vis democratization from below---Svolik's intra-elite bargaining problems can be seen as a further layer of strategic complexity underlying Boix's unitary-elite assumption.
    \item \textbf{Lipset (modernization theory)}: Boix directly tests and largely rejects the pure economic-development account, arguing that it is not income \emph{per se} but the distribution and mobility of assets that drives democratization---development matters primarily insofar as it changes the composition of the asset base (from land to mobile capital).
    \item \textbf{Ostrom 1990}: A more distant connection, but both scholars are centrally concerned with how the distribution of costs and benefits among self-interested actors shapes institutional stability and change; where Ostrom focuses on decentralized, community-level solutions to collective-action problems over common-pool resources, Boix focuses on macro-level class conflict over the distribution of national income via the state.
    \item \textbf{Putnam 1993}: A useful counterpoint in the causal story of institutional performance/emergence: Putnam emphasizes accumulated social capital and civic tradition as the historical driver of institutional quality, while Boix emphasizes material class interest and asset structure---two largely independent (though not mutually exclusive) causal logics for explaining cross-national/cross-regional institutional variation.
    \item \textbf{Cox 1997}: Cox's work on electoral systems and strategic coordination addresses the ``downstream'' question of how democratic institutions, once established, structure political competition; Boix's book addresses the ``upstream'' question of why and when democratic institutions come to exist at all---the two projects can be read as sequential stages of a full theory of democratic politics.
    \item \textbf{Olson 1965}: Boix's model implicitly relies on treating ``the elite'' and ``the poor'' as capable of coherent strategic action despite Olson's collective-action problem; this is a simplifying assumption Boix shares with much of the formal democratization literature, and a point at which Olson's logic could be used to critique or refine the unitary-actor assumptions in Boix's game.
\end{itemize}

\end{document}
